\documentclass[11pt,a4paper]{article}
\usepackage[margin=1in]{geometry}
\usepackage{booktabs}
\usepackage{graphicx}
\usepackage{multirow}
\usepackage{tabularx}
\usepackage{amsmath}
\usepackage{amssymb}
\usepackage{bm}
\usepackage{float}
\usepackage{caption}
\usepackage{subcaption}
\usepackage[utf8]{inputenc}
\usepackage[T1]{fontenc}
\usepackage{hyperref}
\usepackage[numbers]{natbib}
\usepackage{setspace}
\usepackage{hanging}
\graphicspath{{/Users/alberto/Documents/projects/GWP_1/Porfolio-Management/GWP2/}}

\title{MScFE 652 Portfolio Management \\ Group Work Project \#2 }
\author{Seun Oluwafunmilayo Oludayo\\ Megha Shukla\\ Alberto Hernandez-Espinosa}
\date{June 2026}

\begin{document}

\maketitle

\begin{abstract}
This project is the result of development of project MScFE 652 Portfolio Management.
\end{abstract}

\section*{AI Usage Disclosure}
Artificial intelligence tools were used for:
\begin{itemize}
    \item Verifying the accuracy of data sourcing and references;
    \item Validating the mathematical implementation of used algorithms
    \item Assisting in code optimization and debugging to ensure reproducibility of results.
\end{itemize}
All interpretations, conclusions, and original content presented in this report were developed independently by the authors, and all AI-assisted work has been reviewed and verified for accuracy.


\begin{table}[H]
\centering
\caption{Project snapshot}
\label{tab:project_snapshot}
{\small\setlength{\tabcolsep}{4pt}\renewcommand{\arraystretch}{1.15}\begin{tabularx}{\textwidth}{>{\raggedright\arraybackslash}p{4.2cm} X}
\toprule
Item & Value \\
\midrule
Asset universe & 12 assets across equities, gold, treasuries, and REITs (Table~\ref{tab:assets}). \\
Training window & Jan 1, 2024 -- Sep 30, 2025 (438 trading days). \\
Test window & Oct 1, 2025 -- Dec 31, 2025 (63 trading days). \\
Main constraint & Fully invested; no shorts; $0 \le w_i \le 0.20$. \\
Sizing choices & Kelly (100\%), Double-Kelly (200\%), Half-Kelly (50\%), plus Denoised-Kelly, Detoned-Kelly, Denoised+Detoned-Kelly, and HRP in the final ML comparison. \\
Main test result & Double-Kelly returns the most, but HRP has the strongest Sharpe ratio and the smallest drawdown in the final seven-strategy comparison (Table~\ref{tab:step6_results}). \\
Next improvement idea & Cross-validate the Kelly family and compare denoising, clustering, and detoning in one integrated workflow (Steps~3--6). \\
\bottomrule
\end{tabularx}}
\end{table}

\section{Introduction}

This project is about building a simple growth-focused portfolio and then checking how it holds up out-of-sample. Table~\ref{tab:intro_summary} gives the main ideas without a long setup.

\begin{table}[H]
\centering
\caption{Introduction summary (plain language)}
\label{tab:intro_summary}
{\small\setlength{\tabcolsep}{4pt}\renewcommand{\arraystretch}{1.15}\begin{tabularx}{\textwidth}{>{\raggedright\arraybackslash}p{3.8cm} X >{\raggedright\arraybackslash}p{3.2cm}}
\toprule
Topic & What it means here & Reference \\
\midrule
Portfolio choice & Pick weights to balance return and risk in a disciplined way. & Markowitz (1952) \\
Kelly idea & Focus on long-run growth by working with log returns. & Kelly (1956) \\
Practical constraints & Keep weights realistic: fully invested, no shorts, 20\% cap per asset. & Step~1 \\
Why denoising & Covariance estimates can be noisy, and that noise can push weights around. & Laloux et al.\ (1999); L\'{o}pez de Prado (2020) \\
What we report & Build weights in-sample, then compare Kelly/Double/Half out-of-sample. & Steps~1--2 \\
\bottomrule
\end{tabularx}}
\end{table}

\section{Asset Universe and Data}

\subsection{Portfolio Construction}


\begin{table}[H]
\centering
\caption{Portfolio Asset Universe}
\label{tab:assets}
\begin{tabular}{lll}
\toprule
Asset & Name & Sector / Asset Class \\
\midrule
AAPL  & Apple Inc.                    & Technology \\
MSFT  & Microsoft Corp.               & Technology \\
AMZN  & Amazon.com Inc.               & Consumer Discretionary \\
JPM   & JPMorgan Chase \& Co.         & Financials \\
JNJ   & Johnson \& Johnson            & Healthcare \\
XOM   & Exxon Mobil Corp.             & Energy \\
PG    & Procter \& Gamble Co.         & Consumer Staples \\
NEE   & NextEra Energy Inc.           & Utilities \\
BRK-B & Berkshire Hathaway B          & Diversified Conglomerate \\
GLD   & SPDR Gold Shares ETF          & Commodities \\
TLT   & iShares 20+ Year Treasury ETF & Fixed Income \\
VNQ   & Vanguard Real Estate ETF      & Real Estate \\
\bottomrule
\end{tabular}
\end{table}

\begin{table}[H]
\centering
\caption{Data source and return construction}
\label{tab:data_handling}
{\small\setlength{\tabcolsep}{4pt}\renewcommand{\arraystretch}{1.15}\begin{tabularx}{\textwidth}{>{\raggedright\arraybackslash}p{4.0cm} X}
\toprule
Data item & How it was handled \\
\midrule
Prices & Daily adjusted close, downloaded from Yahoo Finance. \\
Returns & Log-returns $r_{i,t} = \ln(P_{i,t}/P_{i,t-1})$ computed for each asset and day. \\
\bottomrule
\end{tabularx}}
\end{table}

\subsection{Data Periods}

\begin{table}[H]
\centering
\caption{Sample windows and setup}
\label{tab:sample_windows}
{\small\setlength{\tabcolsep}{4pt}\renewcommand{\arraystretch}{1.15}\begin{tabularx}{\textwidth}{>{\raggedright\arraybackslash}p{4.2cm} >{\raggedright\arraybackslash}p{6.3cm} >{\raggedright\arraybackslash}p{3.0cm}}
\toprule
Piece & Dates & Size \\
\midrule
In-sample (training) & Jan 1, 2024 -- Sep 30, 2025 & 438 trading days \\
Out-of-sample (testing) & Oct 1, 2025 -- Dec 31, 2025 & 63 trading days \\
Risk-free rate & 4.5\% p.a. (annualised) & Used for excess return \\
\bottomrule
\end{tabularx}}
\end{table}

\subsection{Asset Correlations}

Figure~\ref{fig:corr} shows the correlation matrix over the training
period that are summarised in Table~\ref{tab:key_correlations},
and the broader cross-asset dynamics support a multi-asset allocation strategy:
\begin{table}[H]
\centering
\caption{What correlation say (training period)}
\label{tab:corr_takeaways}
{\small\setlength{\tabcolsep}{4pt}\renewcommand{\arraystretch}{1.15}\begin{tabularx}{\textwidth}{>{\raggedright\arraybackslash}p{4.2cm} X}
\toprule
Pattern & What it suggests \\
\midrule
Tech moves together & Tech names show stronger correlation, so they can behave like one single risk bucket. \\
GLD/TLT hedge role & GLD and TLT are uncorrelated (or slightly negative) with equities, so they can help diversification. \\
Room for diversification & Correlations are not all high, so mixing assets can still reduce portfolio swings. \\
\bottomrule
\end{tabularx}}
\end{table}

\begin{table}[H]
\centering
\caption{Key Pairwise Correlations (Training Period: January 1, 2024 -- September 30, 2025)}
\label{tab:key_correlations}
\begin{tabular}{p{2.2cm}p{5.3cm}c}
\toprule
Pair               & Comment                                      & Correlation \\
\midrule
AAPL-MSFT             & Technology sector peers                              & 0.50 \\
AMZN-MSFT             & Technology/consumer discretionary sector peers        & 0.67 \\
GLD-VNQ             & Gold commodity vs. broad equity market               & $\approx 0$ / $<0$ \\
TLT-VNQ             & Long-duration Treasuries vs. growth sectors       & $\approx 0$ / $<0$ \\
\bottomrule
\end{tabular}
\end{table}

\begin{figure}[H]
\centering
\includegraphics[width=0.85\textwidth]{correlation_heatmap.png}
\caption{Asset correlation matrix over the training period
(January~1,~2024 -- September~30,~2025).
Darker green indicates higher positive correlation.}
\label{fig:corr}
\end{figure}

\section{Step 1: Optimal Kelly Portfolio}

\subsection{Theoretical Framework}

We used a standard Kelly-style mean--variance objective on daily log-returns (Kelly 919--920):

\[
\max_{\bm{w}} \;\; \bm{w}^\top \bm{\mu} - \tfrac{1}{2}\, \bm{w}^\top \bm{\Sigma}\, \bm{w}
\]

Table~\ref{tab:kelly_setup} summarises our setup and the constraints we used.

\begin{table}[H]
\centering
\caption{Kelly optimisation setup (what each piece means)}
\label{tab:kelly_setup}
{\small\setlength{\tabcolsep}{4pt}\renewcommand{\arraystretch}{1.15}\begin{tabularx}{\textwidth}{>{\raggedright\arraybackslash}p{3.2cm} X}
\toprule
Piece & Meaning \\
\midrule
$\bm{w}$ & Portfolio weights across the $N$ assets. \\
$\bm{\mu}$ & Daily excess mean log-returns (net of the risk-free rate), estimated in-sample. \\
$\bm{\Sigma}$ & Daily covariance matrix of log-returns, estimated in-sample. \\
Unconstrained solution & $\bm{w}^* = \bm{\Sigma}^{-1}\bm{\mu}$ (Thorp 389). \\
Constraints used here & Fully invested; no shorting; per-asset cap of 20\% (shown below). \\
Numerical solve & Quadratic program solved with \texttt{cvxpy} (CLARABEL). \\
\bottomrule
\end{tabularx}}
\end{table}

The unconstrained analytical solution is (Thorp 389):

\[
\bm{w}^* = \bm{\Sigma}^{-1} \bm{\mu}
\]

We add practical constraints to keep the weights realistic:

\begin{align*}
    \sum_{i=1}^{N} w_i &= 1 \quad \text{(fully invested)} \\
    0 \leq w_i &\leq 0.20 \quad \forall\, i \quad \text{(no short-selling; 20\% cap per asset)}
\end{align*}

We used
\texttt{cvxpy} convex optimisation library with the CLARABEL solver to solve numerically this quadratic problem.

\subsection{Optimal Weights}

The constrained Kelly optimisation produces the weights shown in
Table~\ref{tab:kelly_weights} and Figure~\ref{fig:weights}.

\begin{table}[H]
\centering
\caption{Constrained Kelly Portfolio Weights ($\leq 20\%$ per asset)}
\label{tab:kelly_weights}
\begin{tabular}{lrr}
\toprule
Asset & Kelly Weight (\%) & Double-Kelly Weight (\%) \\
\midrule
AMZN  & 20.00 & 40.00 \\
BRK-B & 20.00 & 40.00 \\
GLD   & 20.00 & 40.00 \\
JPM   & 20.00 & 40.00 \\
MSFT  & 20.00 & 40.00 \\
AAPL  &  0.00 &  0.00 \\
JNJ   &  0.00 &  0.00 \\
NEE   &  0.00 &  0.00 \\
PG    &  0.00 &  0.00 \\
TLT   &  0.00 &  0.00 \\
VNQ   &  0.00 &  0.00 \\
XOM   &  0.00 &  0.00 \\
\midrule
Total & 100.00 & 200.00 \\
\bottomrule
\end{tabular}
\end{table}

\begin{figure}[H]
\centering
\includegraphics[width=0.90\textwidth]{kelly_weights.png}
\caption{Left: Constrained Kelly weights with the 20\% cap (dashed red line).
Right: Side-by-side comparison of Kelly and Double-Kelly allocations.}
\label{fig:weights}
\end{figure}

\subsection{Interpretation of Weights}

The solution ends up being very ``all or nothing'': a few names at the 20\% cap and the rest at zero. Table~\ref{tab:bang_bang} explains why that happens, and Table~\ref{tab:weight_interpretation} gives the simple rationale by asset group.

\begin{table}[H]
\centering
\caption{Why the weights look ``bang-bang'' in this case}
\label{tab:bang_bang}
{\small\setlength{\tabcolsep}{4pt}\renewcommand{\arraystretch}{1.15}\begin{tabularx}{\textwidth}{>{\raggedright\arraybackslash}p{4.4cm} X}
\toprule
What we see & Simple reason \\
\midrule
Five assets at 20\% & The unconstrained Kelly solution wants to concentrate on the strongest contributors, so once the cap binds it fills up the top picks first. \\
Many assets at 0\% & Some assets do not improve the risk/return trade-off enough after correlations are considered, so the optimiser leaves them out under the cap. \\
Clustering at the bounds & With a hard cap and no shorting, the optimiser often ends up pushing weights to the extremes (0\% or 20\%) instead of spreading them evenly. \\
\bottomrule
\end{tabularx}}
\end{table}
\begin{table}[H]
\centering
\caption{Rationale for Portfolio Weight Allocations}
\label{tab:weight_interpretation}
{\setlength{\tabcolsep}{4pt}\begin{tabularx}{\textwidth}{l l X}
\toprule
Group & Tickers & Explanation\\
\midrule
\multirow{5}{*}{Included (20\% weight)} 
& AMZN, MSFT & Did well over the period, especially compared with the amount of day-to-day swings \\
& JPM & Benefited from the higher interest-rate backdrop and held up better than most \\
& BRK-B & Generally steadier than the rest and still delivered solid returns \\
& GLD & Helped balance the equity risk since it did not move in step with the stocks \\
\midrule
\multirow{2}{*}{Excluded (0\% weight)}
& JNJ, NEE, PG, XOM, VNQ & Nothing wrong with them, but they did not add enough return or diversification in this setup \\
& TLT & Did not add much extra protection on top of GLD, and long bonds were hurt as rates moved up in 2024--2025 \\
\bottomrule
\end{tabularx}}
\end{table}

\section{Step 2: Out-of-Sample (test dataset) Performance Comparison}

\subsection{Strategy and Performance Metrics Overview}

Tables~\ref{tab:strategy_defs} and~\ref{tab:metric_defs} define the three sizing choices and the metrics used for comparison.

\begin{table}[H]
\centering
\caption{Out-of-Sample Strategy Definitions}
\label{tab:strategy_defs}
\begin{tabularx}{\textwidth}{l l X}
\toprule
Strategy & Gross Notional Exposure & Description \\
\midrule
Kelly        & 100\% & Just the basic constrained allocation using the raw weights from Table~\ref{tab:kelly_weights}, all money put in with no borrowing to add extra positions. \\
Double-Kelly & 200\% & Levered version where every Kelly weight is doubled, so the portfolio is 2:1 on gross exposure. \\
Half-Kelly   & 50\%  & More cautious version where every Kelly weight is cut in half and the rest is held in the risk-free asset. \\
\bottomrule
\end{tabularx}
\end{table}


\begin{table}[H]
\centering
\caption{Performance Metric Definitions}
\label{tab:metric_defs}
{\small\setlength{\tabcolsep}{4pt}\renewcommand{\arraystretch}{1.15}\begin{tabularx}{\textwidth}{>{\raggedright\arraybackslash}p{3.3cm} >{\raggedright\arraybackslash}p{5.2cm} X}
\toprule
Metric & Formula & Interpretation \\
\midrule
Cumulative Return & \(\displaystyle R_{cum} = e^{\sum_t r_t} - 1\) & How much the whole portfolio grew (or shrank) in percentage terms over the whole test window, figured by adding up all the daily log returns we tracked. \\
Sharpe Ratio & \(\displaystyle S = \frac{\bar{r}_{excess}}{\sigma_r}\sqrt{252}\) & Annual number that shows if returns were worth the risk taken, compares average extra returns above the 4.5\% risk-free rate to how much daily returns bounced around, then scaled to a full year. \\
Maximum Drawdown (MDD) & \(\displaystyle \text{MDD} = \min_t \frac{V_t - M_t}{M_t}, \;\; M_t=\max_{s \leq t} V_s\) & Biggest single drop from a high point to a low point the portfolio saw during the test, basically the worst loss anyone holding this would've had to sit through at any point. \\
95\%-CVaR (daily) & Empirical mean of the worst 5\% of daily returns & Average loss on the really bad trading days, the 5\% worst ones, gives a sense of how bad those terrible days usually were, not just the single worst one. \\
\bottomrule
\end{tabularx}}
\end{table}
\subsection{Results}

\begin{table}[H]
\centering
\caption{Out-of-sample performance -- October 1 -- December 31, 2025}
\label{tab:oos_results}
\begin{tabular}{lcccc}
\toprule
Strategy & Cumulative Return (\%) & Sharpe Ratio & Max Drawdown (\%) & 95\%-CVaR daily (\%) \\
\midrule
Kelly        &  2.39 & 0.40 &  $-5.54$ & $-1.44$ \\
Double-Kelly &  4.83 & 0.59 & $-10.77$ & $-2.88$ \\
Half-Kelly   &  1.19 & 0.02 &  $-2.81$ & $-0.72$ \\
\bottomrule
\multicolumn{5}{l}{\small Note: CVaR and MDD are negative values (losses). Risk-free rate: 4.5\% p.a.}
\end{tabular}
\end{table}

\begin{figure}[H]
\centering
\includegraphics[width=\textwidth]{cumulative_returns_oos.png}
\caption{Cumulative returns of the three strategies -- out-of-sample period (October~1 -- December~31, 2025).}
\label{fig:cumret}
\end{figure}

\begin{figure}[H]
\centering
\includegraphics[width=\textwidth]{drawdown_oos.png}
\caption{Portfolio drawdown paths -- out-of-sample period. The shaded area highlights the Double-Kelly drawdown, which peaked at $-10.77\%$ around mid-November 2025.}
\label{fig:drawdown}
\end{figure}

\begin{figure}[H]
\centering
\includegraphics[width=\textwidth]{metrics_comparison.png}
\caption{Summary of the four risk and performance metrics for each strategy.}
\label{fig:metrics}
\end{figure}

\begin{figure}[H]
\centering
\includegraphics[width=\textwidth]{return_distributions.png}
\caption{Empirical daily return distributions and 5th percentile (VaR) thresholds for each strategy -- out-of-sample period.}
\label{fig:dist}
\end{figure}

\subsection{Analysis and Discussion}

The results in Table~\ref{tab:oos_results} are easier to read when we translate the
numbers into a few simple takeaways. Table~\ref{tab:oos_takeaways} summarises the
main points, and Table~\ref{tab:oos_reco} gives a practical interpretation.

\begin{table}[H]
\centering
\caption{Out-of-sample results: quick interpretation}
\label{tab:oos_takeaways}
{\small\setlength{\tabcolsep}{4pt}\renewcommand{\arraystretch}{1.15}\begin{tabularx}{\textwidth}{>{\raggedright\arraybackslash}p{3.0cm} >{\raggedright\arraybackslash}p{2.6cm} >{\raggedright\arraybackslash}p{2.8cm} >{\raggedright\arraybackslash}p{2.6cm} X}
\toprule
Metric & Kelly & Double-Kelly & Half-Kelly & Explanation \\
\midrule
Cumulative return & 2.39\% & 4.83\% & 1.19\% & Double-Kelly is roughly ``twice the exposure, roughly twice the result'' over this short window (Thorp 395--396). \\
Sharpe ratio & 0.40 & 0.59 & 0.02 & In practice, the risk-free adjustment matters: when returns are small, subtracting the risk-free rate can noticeably change the ratio, especially over a short test period. \\
Max drawdown & $-5.54\%$ & $-10.77\%$ & $-2.81\%$ & Losses scale up with leverage too. The worst dip for Double-Kelly is about twice Kelly, and it happened around mid-November 2025 during a broad equity sell-off. \\
95\%-CVaR (daily) & $-1.44\%$ & $-2.88\%$ & $-0.72\%$ & On the worst 5\% of days, Double-Kelly loses about twice as much as Kelly, which is the key cost of leverage. \\
\bottomrule
\end{tabularx}}
\end{table}

\begin{table}[H]
\centering
\caption{Practical takeaways from the three leverage choices}
\label{tab:oos_reco}
{\small\setlength{\tabcolsep}{4pt}\renewcommand{\arraystretch}{1.15}\begin{tabularx}{\textwidth}{>{\raggedright\arraybackslash}p{3.6cm} >{\raggedright\arraybackslash}p{3.0cm} X}
\toprule
Investor style & Suggested strategy & Brief reason \\
\midrule
Balanced risk appetite & Kelly & Keeps decent upside (2.39\%) while drawdowns stay more manageable than the leveraged version. \\
High risk tolerance and leverage capacity & Double-Kelly & Higher return, but you must be comfortable with larger drawdowns and worse tail losses. \\
More cautious sizing / uncertainty is high & Half-Kelly & Smaller drawdowns, but risk-adjusted performance can look weak over short windows once the risk-free rate is accounted for (L\'{o}pez de Prado 42--44). \\
\bottomrule
\end{tabularx}}
\end{table}

\section{Step 3: K-fold Cross-Validation Backtesting}

\subsection{Walk-forward K-fold design}

Instead of judging the portfolio rules on one
out-of-sample window only, we rotate the test window across the full 2024--2025 sample and
re-estimate the model each time. The standard random K-fold procedure is not appropriate for
financial data because shuffling would let future observations leak into the training set. We used a walk-forward variant in which the training window
always ends before the test window begins, and we add a 5-day embargo at the boundary to
reduce contamination from short-run autocorrelation.

\begin{table}[H]
\centering
\caption{Step 3 backtesting design}
\label{tab:step3_design}
{\small\setlength{\tabcolsep}{4pt}\renewcommand{\arraystretch}{1.15}\begin{tabularx}{\textwidth}{>{\raggedright\arraybackslash}p{4.2cm} X}
\toprule
Item & Implementation in the script \\
\midrule
Dataset used & Full sample from Jan 3, 2024 to Dec 31, 2025 (501 daily log-return observations after differencing). \\
Number of folds & 5 walk-forward folds. \\
Initial training size & 86 observations. \\
Test size per fold & 83 observations. \\
Embargo & 5 trading days between training end and test start. \\
Strategies compared & Kelly, Double-Kelly, and Half-Kelly. \\
Metrics reported & Cumulative return, Sharpe ratio, Sortino ratio, Maximum Drawdown, and 95\%-CVaR. \\
\bottomrule
\end{tabularx}}
\end{table}

\subsection{Cross-validation results}

Double-Kelly has the
highest mean return and the highest mean Sharpe ratio across folds, but it is also the least stable.
Its standard deviation across folds is almost as large as its mean Sharpe ratio, and its average
drawdown and tail loss are materially worse than the alternatives. Kelly sits in the middle, while
Half-Kelly gives the most conservative profile, with smaller losses but also weaker risk-adjusted
payoff.

\begin{table}[H]
\centering
\caption{Step 3 cross-validation summary: mean (standard deviation) across 5 folds}
\label{tab:step3_summary}
{\small\setlength{\tabcolsep}{4pt}\renewcommand{\arraystretch}{1.15}\begin{tabular}{lccccc}
\toprule
Strategy & Cum. Return (\%) & Sharpe & Sortino & Max Drawdown (\%) & 95\%-CVaR (\%) \\
\midrule
Kelly & 4.42 (3.19) & 0.83 (0.95) & 0.78 (0.91) & $-7.44$ (4.73) & $-1.92$ (1.01) \\
Double-Kelly & 9.12 (6.68) & 1.01 (0.99) & 0.95 (0.95) & $-14.15$ (8.50) & $-3.84$ (2.02) \\
Half-Kelly & 2.18 (1.56) & 0.46 (0.87) & 0.43 (0.81) & $-3.82$ (2.49) & $-0.96$ (0.50) \\
\bottomrule
\multicolumn{6}{l}{\small Note: values are mean (standard deviation) across the 5 walk-forward folds.}
\end{tabular}}
\end{table}

\begin{figure}[H]
\centering
\includegraphics[width=0.92\textwidth]{step3_cv_summary.png}
\caption{Step 3 cross-validation summary. Left: mean Sharpe ratio across folds. Right: mean maximum drawdown across folds.}
\label{fig:step3_cv}
\end{figure}

\subsection{Interpretation}

This pattern is close to the logic in previous steps. Leverage can help when the
estimated edge is right, so Double-Kelly often posts the highest fold return, but the same leverage
also magnifies estimation error and tail losses. In practical terms, the walk-forward backtest says
that Double-Kelly is the most aggressive choice, Kelly is the middle-ground choice, and
Half-Kelly is the most robust if parameter uncertainty is the main concern.

\section{Step 4: Improvements Using Denoising, Clustering, and Detoning}

\subsection{The Estimation Error Problem in Covariance Matrices}

Here we estimated the covariance matrix $\bm{\Sigma}$, and found that it is noisy. Estimates are noisy, and that noise can move weights around. Table~\ref{tab:estimation_noise} puts the idea in the context of our data.  

\begin{table}[H]
\centering
\caption{Estimation noise: what it means for this project}
\label{tab:estimation_noise}
{\small\setlength{\tabcolsep}{4pt}\renewcommand{\arraystretch}{1.15}\begin{tabularx}{\textwidth}{>{\raggedright\arraybackslash}p{4.2cm} X}
\toprule
Piece & Plain meaning \\
\midrule
What is being estimated & Expected returns $\bm{\mu}$ and covariance $\bm{\Sigma}$ from historical data. \\
Why it is a problem & Random variation can look like a pattern, and the optimiser may chase it. \\
Why covariances matter a lot & The Kelly-style solution depends heavily on $\bm{\Sigma}^{-1}$, which can amplify estimation noise. \\
Why this study still makes sense & With $N=12$ and $T=438$, estimates are workable but still noisy enough that cleaning $\bm{\Sigma}$ can help (Laloux et al.\ 1467). \\
\bottomrule
\end{tabularx}}
\end{table}

\subsection{Random Matrix Theory and the Marc\u{e}nko-Pastur Distribution}

RMT gives a rule of thumb for which eigenvalues look like noise, based on $q=N/T$ (L\'{o}pez de Prado 58; Laloux et al.\ 1467--1468).

\[
\lambda^{\pm} = \sigma^2 \left(1 \pm \frac{1}{\sqrt{q}}\right)^2
\]

\begin{table}[H]
\centering
\caption{RMT quantities used in this study}
\label{tab:rmt_params}
{\small\setlength{\tabcolsep}{4pt}\renewcommand{\arraystretch}{1.15}\begin{tabularx}{\textwidth}{>{\raggedright\arraybackslash}p{4.0cm} X}
\toprule
Quantity & How we use it \\
\midrule
$q = N/T$ & Ratio that controls the noise band for eigenvalues. Here $q = 12/438 \approx 0.0274$. \\
$\lambda^+$ & Upper bound of the ``mostly noise'' region. Here $\lambda^+ \approx 1.37$. \\
Interpretation & Eigenvalues at/below $\lambda^+$ are treated as noise-like; bigger eigenvalues are treated as structure worth keeping. \\
\bottomrule
\end{tabularx}}
\end{table}

\subsection{The Denoising Procedure}

Table~\ref{tab:denoise_steps} summarises the standard Marc\u{e}nko-Pastur denoising
recipe (L\'{o}pez de Prado 60--63).

\begin{table}[H]
\centering
\caption{Denoising procedure (Marc\u{e}nko-Pastur): plain summary}
\label{tab:denoise_steps}
{\small\setlength{\tabcolsep}{4pt}\renewcommand{\arraystretch}{1.15}\begin{tabularx}{\textwidth}{>{\raggedright\arraybackslash}p{2.0cm} X}
\toprule
Step & What you do \\
\midrule
1 & Build the sample correlation matrix $\hat{\bm{C}}$ from log-returns. \\
2 & Compute eigenvalues and eigenvectors: $\hat{\bm{C}} = \bm{V}\bm{\Lambda}\bm{V}^\top$. \\
3 & Use $\lambda^+$ to decide which eigenvalues look like noise (those at or below $\lambda^+$). \\
4 & Replace the ``noise'' eigenvalues with a single average level (chosen so the total variance is preserved), then rebuild a denoised $\tilde{\bm{C}}$. \\
5 & Convert back to a covariance matrix using the original asset standard deviations. \\
\bottomrule
\end{tabularx}}
\end{table}

\subsection{Benefits for Kelly Optimisation}

Applying denoising before solving the constrained Kelly problem is mainly about making
the inputs less fragile. Table~\ref{tab:denoise_benefits} summarises the practical
benefits.

\begin{table}[H]
\centering
\caption{Why denoising can help the Kelly optimisation}
\label{tab:denoise_benefits}
{\small\setlength{\tabcolsep}{4pt}\renewcommand{\arraystretch}{1.15}\begin{tabularx}{\textwidth}{>{\raggedright\arraybackslash}p{3.4cm} X}
\toprule
Benefit & Plain interpretation \\
\midrule
Less overfitting in $\bm{\Sigma}^{-1}$ & The optimiser can overreact to small, noisy eigenvalues because inversion amplifies them. Denoising reduces this effect and usually leads to more stable weights (L\'{o}pez de Prado 64--65). \\
Better numerical stability & A better-conditioned covariance matrix makes optimisation and inversion less sensitive to small data changes. \\
Less ``whipsaw'' turnover & Weights tend to move less from one rebalance to the next, which matters once you consider transaction costs. \\
Empirical support & Denoised covariances often improve out-of-sample Sharpe ratios in practice, especially when $T/N$ is not very large (L\'{o}pez de Prado 66). \\
\bottomrule
\end{tabularx}}
\end{table}

\subsection{Application to Our Portfolio}

If we denoise the covariance estimates for this 12-asset universe, the main expectation is
that the allocation would be less ``all-in'' on the same five names at the 20\% cap.
In simple terms, denoising can make the optimiser less confident in very concentrated bets.
Table~\ref{tab:denoise_impact} states the likely direction of change.

\begin{table}[H]
\centering
\caption{Expected effect of denoising in this study (directional)}
\label{tab:denoise_impact}
{\small\setlength{\tabcolsep}{4pt}\renewcommand{\arraystretch}{1.15}\begin{tabularx}{\textwidth}{>{\raggedright\arraybackslash}p{4.0cm} X}
\toprule
What would change & What we expect to see \\
\midrule
Weight concentration & A less extreme ``five-at-the-cap'' pattern, with some weight spreading across more assets. \\
Previously excluded assets (e.g., TLT, VNQ) & They may look more helpful once noisy covariance effects are reduced, so they could receive small positive weights in the constrained solution. \\
Tail risk measures (MDD, CVaR) & Likely lower, since the portfolio may be less concentrated in correlated equity exposures. \\
Raw cumulative return & Possibly a bit lower, since reducing concentration can give up some upside when the concentrated names keep performing. \\
\bottomrule
\end{tabularx}}
\end{table}

Steps~5 and~6 apply this and compare the denoised version out-of-sample.

\subsection{Improvements Using Clustering: Hierarchical Risk Parity}

Here was used clustering from a practical problem in mean--variance style optimisation:
if inputs are noisy, the optimiser can mistake noise for signal and over-allocate to whatever looks
best in-sample. HRP avoids that fragility by side-stepping expected return estimation and by
allocating risk across clusters rather than pushing directly through a covariance-matrix inverse.

\[
d(i,j)=\sqrt{\frac{1-\rho_{ij}}{2}}
\]

This distance turns correlations into a clustering input. Assets that move similarly end up closer
together, and hierarchical clustering can then group them before the recursive allocation step.
The final weights are set by recursive bisection so that no single cluster dominates total risk.

\begin{table}[H]
\centering
\caption{HRP in plain language}
\label{tab:hrp_summary}
{\small\setlength{\tabcolsep}{4pt}\renewcommand{\arraystretch}{1.15}\begin{tabularx}{\textwidth}{>{\raggedright\arraybackslash}p{3.0cm} X}
\toprule
Stage & What happens \\
\midrule
Distance step & Convert the correlation matrix into pairwise distances using \(d(i,j)=\sqrt{(1-\rho_{ij})/2}\). \\
Clustering step & Build a hierarchical tree so that similar assets sit in the same branch. \\
Quasi-diagonalisation & Reorder the covariance matrix so related assets sit next to each other. \\
Recursive bisection & Split the asset tree into sub-clusters and allocate weight inversely to cluster variance. \\
Practical effect & The portfolio spreads risk across groups instead of chasing the highest estimated mean return. \\
\bottomrule
\end{tabularx}}
\end{table}

The resulting HRP weights are much more diversified than the constrained Kelly solution. The
largest allocations are TLT (20.82\%), GLD (15.66\%), PG (10.22\%), VNQ (9.40\%), JNJ
(9.39\%), and XOM (8.71\%). That is exactly the kind of pattern clustering is designed to
produce: less concentration in a few equity names and more balance across different risk buckets.

\subsection{Improvements Using Detoning}

Even after denoising, the first
eigenvector in the correlation matrix can still represent a broad market mode that pushes many
assets in the same direction. If the goal is diversification, that common factor can make the
optimiser look more confident than it should be about concentrated bets. Detoning removes that
dominant mode and then re-normalises the matrix.

\[
\bm{C}_{\text{detoned}} = \bm{C} - \lambda_{\text{mkt}} \bm{v}_{\text{mkt}} \bm{v}_{\text{mkt}}^\top
\]

\begin{table}[H]
\centering
\caption{What detoning changes}
\label{tab:detone_summary}
{\small\setlength{\tabcolsep}{4pt}\renewcommand{\arraystretch}{1.15}\begin{tabularx}{\textwidth}{>{\raggedright\arraybackslash}p{3.6cm} X}
\toprule
Question & Practical answer \\
\midrule
What is removed? & The first eigenvector, interpreted as the broad market factor. \\
Why remove it? & It can dominate the covariance structure and hide genuine diversification opportunities. \\
What happened in our run? & The largest eigenvalue was 3.5170, equal to 29.3\% of total trace before detoning. \\
What happens after denoising first? & The stacked denoised$\rightarrow$detoned matrix still has a dominant market factor, with eigenvalue 3.5243 (29.4\% of trace), which is then removed. \\
Expected portfolio effect & Slightly less concentration and a better chance that diversifying assets receive positive weight. \\
\bottomrule
\end{tabularx}}
\end{table}

\section{Step 5: Applying the ML Improvements}

\subsection{What was implemented}
We applied the following improvements to the
same 12-asset universe used in Step~1:

\begin{table}[H]
\centering
\caption{Step 5 methodology menu}
\label{tab:step5_methods}
{\small\setlength{\tabcolsep}{4pt}\renewcommand{\arraystretch}{1.15}\begin{tabularx}{\textwidth}{>{\raggedright\arraybackslash}p{3.2cm} >{\raggedright\arraybackslash}p{3.9cm} X}
\toprule
Tool & Core merit & Main drawback \\
\midrule
HRP & Stable, cluster-based diversification without matrix inversion. & Does not explicitly maximise growth through expected returns. \\
Constrained Kelly & Targets long-run compounding growth while keeping a 20\% cap and no shorting. & Still sensitive to estimation error in the inputs. \\
Double-Kelly & Pushes the growth logic further through leverage. & Can magnify drawdowns and tail losses very quickly. \\
Denoising & Cleans the covariance estimate by shrinking noise-like eigenvalues. & By itself, it may not change the final portfolio if the original cap constraints already bind. \\
Detoning & Removes the broad market mode and can open room for more diversification. & Can lower concentration but does not guarantee better raw return. \\
\bottomrule
\end{tabularx}}
\end{table}

\subsection{Explaining results}

\begin{table}[H]
\centering
\caption{Step 5 log excerpt reference}
\label{tab:step5_logs_ref}
{\small\setlength{\tabcolsep}{4pt}\renewcommand{\arraystretch}{1.15}\begin{tabularx}{\textwidth}{>{\raggedright\arraybackslash}p{4.0cm} X}
\toprule
Logged item & Result from the script run \\
\midrule
Denoising noise band & \([0.6960,\ 1.3589]\) \\
Noisy eigenvalues & 4 of 12 \\
Market eigenvalue before detoning & 3.5170 (29.3\% of trace) \\
Market eigenvalue after denoising, before stacked detoning & 3.5243 (29.4\% of trace) \\
\bottomrule
\end{tabularx}}
\end{table}

\begin{verbatim}
=== Denoising ===
  MP noise band       : [0.6960,  1.3589]
  Noisy eigenvalues   : 4 / 12
Denoised covariance matrix computed.

=== Detoning ===
  Market eigenvalue   : 3.5170  (29.3% of trace)

=== Denoised -> Detoned (stacked) ===
  Market eigenvalue   : 3.5243  (29.4% of trace)
\end{verbatim}

\subsection{Portfolio weights after the improvements}

The denoising-only portfolio is practically identical to the baseline constrained Kelly solution:
the same five assets remain pinned at the 20\% cap. This is a useful result in itself, because it
shows that the binding cap is already strong enough that a modest covariance cleaning step does
not move the optimum. Detoning changes the solution slightly by allowing a small positive weight
to AAPL and reducing MSFT below the cap. The stacked denoised+detoned version pushes this
effect a bit further. HRP is the most visibly different portfolio because it spreads risk across all
12 assets.

\begin{table}[H]
\centering
\caption{Step 5 portfolio weights (\%)}
\label{tab:step5_weights}
{\scriptsize\setlength{\tabcolsep}{4pt}\renewcommand{\arraystretch}{1.12}\begin{tabular}{lrrrrr}
\toprule
Ticker & Kelly & Denoised & Detoned & Denoised+Detoned & HRP \\
\midrule
AAPL  & 0.00 & 0.00 & 0.76 & 1.60 & 4.68 \\
AMZN  & 20.00 & 20.00 & 20.00 & 20.00 & 2.94 \\
BRK-B & 20.00 & 20.00 & 20.00 & 20.00 & 6.04 \\
GLD   & 20.00 & 20.00 & 20.00 & 20.00 & 15.66 \\
JNJ   & 0.00 & 0.00 & 0.00 & 0.00 & 9.39 \\
JPM   & 20.00 & 20.00 & 20.00 & 20.00 & 2.98 \\
MSFT  & 20.00 & 20.00 & 19.25 & 18.40 & 5.47 \\
NEE   & 0.00 & 0.00 & 0.00 & 0.00 & 3.68 \\
PG    & 0.00 & 0.00 & 0.00 & 0.00 & 10.22 \\
TLT   & 0.00 & 0.00 & 0.00 & 0.00 & 20.82 \\
VNQ   & 0.00 & 0.00 & 0.00 & 0.00 & 9.40 \\
XOM   & 0.00 & 0.00 & 0.00 & 0.00 & 8.71 \\
\bottomrule
\end{tabular}}
\end{table}

\begin{figure}[H]
\centering
\includegraphics[width=\textwidth]{step5_weights_comparison.png}
\caption{Step 5 weights comparison across the baseline Kelly portfolio and the ML-enhanced alternatives.}
\label{fig:step5_weights}
\end{figure}

\subsection{Merits and sequence}

We believe that these methods are better seen as complementary tools than
as substitutes. That logic still holds after integration. A practical sequence is:

\begin{table}[H]
\centering
\caption{Suggested sequence for the three improvements}
\label{tab:step5_sequence}
{\small\setlength{\tabcolsep}{4pt}\renewcommand{\arraystretch}{1.15}\begin{tabularx}{\textwidth}{>{\raggedright\arraybackslash}p{1.3cm} X}
\toprule
Step & Reason \\
\midrule
1 & Denoise first, because it removes the most obvious covariance noise before any further transformation. \\
2 & Detone next, because removing the market factor is cleaner once the raw noise has already been reduced. \\
3 & Re-optimise Kelly on the cleaned matrix and benchmark the result against HRP and the original Kelly portfolio. \\
4 & Use HRP as the diversification benchmark, especially if stability and drawdown control matter more than pure growth maximisation. \\
\bottomrule
\end{tabularx}}
\end{table}

\section{Step 6: Do Multiple Improvements Outperform?}

\subsection{Seven-strategy out-of-sample comparison}

Here we compare the original three Kelly sizing rules with four ML-enhanced alternatives over the
same out-of-sample window, October~1 to December~31, 2025. We
compare cumulative return, annualised Sharpe ratio, Sortino ratio, Maximum Drawdown, and
95\%-CVaR.

\begin{table}[H]
\centering
\caption{Step 6 out-of-sample results: all seven strategies}
\label{tab:step6_results}
{\small\setlength{\tabcolsep}{4pt}\renewcommand{\arraystretch}{1.12}\begin{tabular}{lccccc}
\toprule
Strategy & Cum. Return (\%) & Sharpe & Sortino & Max Drawdown (\%) & 95\%-CVaR (\%) \\
\midrule
Kelly & 2.3877 & 0.3995 & 0.3948 & $-5.5364$ & $-1.4394$ \\
Double-Kelly & 4.8324 & 0.5872 & 0.5858 & $-10.7663$ & $-2.8787$ \\
Half-Kelly & 1.1868 & 0.0243 & 0.0239 & $-2.8076$ & $-0.7197$ \\
Denoised-Kelly & 2.3877 & 0.3995 & 0.3948 & $-5.5364$ & $-1.4394$ \\
Detoned-Kelly & 2.4907 & 0.4331 & 0.4287 & $-5.5098$ & $-1.4341$ \\
Denoised+Detoned-Kelly & 2.6064 & 0.4708 & 0.4668 & $-5.4797$ & $-1.4281$ \\
HRP & 2.8174 & 0.9033 & 0.8875 & $-2.5599$ & $-0.7931$ \\
\bottomrule
\end{tabular}}
\end{table}

\begin{figure}[H]
\centering
\includegraphics[width=\textwidth]{step6_cumulative_returns.png}
\caption{Step 6 cumulative returns for all seven strategies over the out-of-sample period.}
\label{fig:step6_cum}
\end{figure}

\begin{figure}[H]
\centering
\includegraphics[width=\textwidth]{step6_metrics_bar.png}
\caption{Step 6 comparison of cumulative return, annualised Sharpe ratio, and maximum drawdown across all seven strategies.}
\label{fig:step6_metrics}
\end{figure}

\subsection{Interpretation of results}

\begin{table}[H]
\centering
\caption{Step 6 interpretation}
\label{tab:step6_takeaways}
{\small\setlength{\tabcolsep}{4pt}\renewcommand{\arraystretch}{1.15}\begin{tabularx}{\textwidth}{>{\raggedright\arraybackslash}p{3.6cm} X}
\toprule
Finding & Interpretation \\
\midrule
Highest raw return & Double-Kelly remains the return leader at 4.83\%, but that gain comes with the worst drawdown and the worst tail loss. \\
Best risk-adjusted result & HRP has the highest Sharpe ratio (0.9033), the highest Sortino ratio (0.8875), and the lowest maximum drawdown ($-2.5599\%$). \\
Best Kelly-family improvement & Denoised+Detoned-Kelly is the strongest Kelly variant once ML adjustments are allowed, improving return and both risk-adjusted ratios relative to baseline Kelly. \\
Denoising alone & Denoised-Kelly is numerically the same as baseline Kelly in this sample because the weight cap still binds the same five assets. \\
Half-Kelly & Half-Kelly cuts risk materially, but the risk-free adjustment leaves its Sharpe ratio close to zero over this short test window. \\
\bottomrule
\end{tabularx}}
\end{table}

As initial conclusion we think think that multiple improvements can help, but the value
depends on what the investor is trying to improve. If the target is the highest raw return,
Double-Kelly still wins. If the target is a better overall balance between return, drawdown, and
tail risk, HRP is the clearest winner. If the investor wants to stay within the Kelly framework,
denoising and detoning work better together than separately.
\section{Conclusion}

\begin{table}[H]
\centering
\caption{Conclusion summary}
\label{tab:conclusion_summary}
{\small\setlength{\tabcolsep}{4pt}\renewcommand{\arraystretch}{1.15}\begin{tabularx}{\textwidth}{>{\raggedright\arraybackslash}p{4.0cm} X}
\toprule
Point & Summary \\
\midrule
In-sample baseline result & With a 20\% cap and no shorts, the original Kelly solution still concentrates in five assets at the cap. \\
Step 3 message & Walk-forward K-fold testing confirms that leverage lifts average return but also makes results much less stable across folds. \\
Step 4 message & Denoising, clustering, and detoning all try to reduce fragility, but they do so in different ways: cleaning noise, spreading risk across clusters, or removing the market mode. \\
Best ML-enhanced Kelly result & Denoised+Detoned-Kelly is the strongest Kelly-family variant in the Q4~2025 out-of-sample test. \\
Best overall trade-off & HRP gives the best overall risk-adjusted performance in the final seven-strategy comparison. \\
\bottomrule
\end{tabularx}}
\end{table}

The traditional Kelly framework is
useful, but its aggressive versions are highly sensitive to estimation error. Cross-validation shows
that leverage can pay when the signal is right, yet the losses scale quickly when it is wrong.
Among the ML-style improvements, denoising and detoning together improve the Kelly family
modestly, while HRP changes the portfolio structure more substantially and delivers the best
balance between return and risk in the final out-of-sample comparison.

\bibliographystyle{plainnat}
\begin{thebibliography}{99}

\bibitem{kelly1956}
Kelly, J. L. ``A New Interpretation of Information Rate.''
\textit{Bell System Technical Journal}, vol.~35, no.~4, 1956, pp.~917--926.

\bibitem{laloux1999}
Laloux, L., et al. ``Noise Dressing of Financial Correlation Matrices.''
\textit{Physical Review Letters}, vol.~83, no.~7, 1999, pp.~1467--1470.

\bibitem{lopezdeprado2016}
L\'{o}pez de Prado, Marcos.
``Building Diversified Portfolios that Outperform Out of Sample.''
\textit{The Journal of Portfolio Management}, vol.~42, no.~4, 2016, pp.~59--69.

\bibitem{lopezdeprado2018}
L\'{o}pez de Prado, Marcos.
\textit{Advances in Financial Machine Learning}.
Wiley, 2018.

\bibitem{lopezdeprado2020}
L\'{o}pez de Prado, Marcos.
\textit{Machine Learning for Asset Managers}.
Cambridge University Press, 2020.

\bibitem{maclean2011}
MacLean, Leonard C., Edward O. Thorp, and William T. Ziemba.
\textit{The Kelly Capital Growth Investment Criterion}.
World Scientific, 2011.

\bibitem{markowitz1952}
Markowitz, Harry. ``Portfolio Selection.''
\textit{The Journal of Finance}, vol.~7, no.~1, 1952, pp.~77--91.

\bibitem{merton1980}
Merton, Robert C.
``On Estimating the Expected Return on the Market: An Exploratory Investigation.''
\textit{Journal of Financial Economics}, vol.~8, no.~4, 1980, pp.~323--361.

\bibitem{michaud1989}
Michaud, Richard O.
``The Markowitz Optimization Enigma: Is ``Optimized'' Optimal?''
\textit{Financial Analysts Journal}, vol.~45, no.~1, 1989, pp.~31--42.

\bibitem{demiguel2009}
DeMiguel, Victor, Lorenzo Garlappi, and Raman Uppal.
``Optimal versus Naive Diversification: How Inefficient is the 1/N Portfolio Strategy?''
\textit{Review of Financial Studies}, vol.~22, no.~5, 2009, pp.~1915--1953.

\bibitem{bailey2014}
Bailey, David H., et al.
``The Probability of Backtest Overfitting.''
\textit{Journal of Computational Finance}, vol.~20, no.~4, 2017, pp.~39--69.

\bibitem{thorp2006}
Thorp, Edward O. ``The Kelly Criterion in Blackjack, Sports Betting, and the
Stock Market.''
\textit{Handbook of Asset and Liability Management}, vol.~1,
edited by S. A. Zenios and W. Ziemba,
North-Holland, 2006, pp.~385--428.

\bibitem{ward1963}
Ward, Joe H.
``Hierarchical Grouping to Optimize an Objective Function.''
\textit{Journal of the American Statistical Association}, vol.~58, no.~301, 1963, pp.~236--244.

\end{thebibliography}

\end{document}
